% FILE: 02_lineage_and_context.tex
% PROJECT: otel-weaver-exp-001-custom-pi-registry
% THESIS ROLE: otel-semconv-evidence-surface-experiment

\section{Lineage and Context}
\label{sec:otel-weaver-exp-001-custom-pi-registry-lineage}

\subsection{2016 Language-Model Lesson}
\label{sec:otel-weaver-exp-001-custom-pi-registry-lineage-2016}

The 2016 language-model experiment in the dissertation's lineage established that local text
imitation does not conserve consequence: a model that reproduces the statistical surface of a
corpus does not thereby reproduce the causal structure that produced that corpus. Applied to
telemetry, the lesson is structurally identical. An OTel span that reproduces the surface
appearance of a process event --- a timestamp, a service name, a status code --- does not thereby
reproduce the process-evidence structure required for conformance replay. The span is an imitation
of an event, not a lawful encoding of it.

EXP-001 is a direct response to this lesson. By requiring that admissible spans carry
\texttt{process.pi.witness.hash} (a BLAKE3 64-character hexadecimal seal),
\texttt{process.pi.token.state\_before} and \texttt{state\_after} (JSON-serialized Petri net
markings), and \texttt{process.pi.lifecycle} (one of: \texttt{schedule}, \texttt{start},
\texttt{suspend}, \texttt{resume}, \texttt{complete}, \texttt{abort}), the registry enforces that
the span carries not merely the surface of a process event but the causal state required to
reconstruct the event in a conformance court. The schema boundary is the technical implementation
of the 2016 lesson: consequence must be conserved at the feedstock layer or not at all.

\subsection{Enterprise Process Gap}
\label{sec:otel-weaver-exp-001-custom-pi-registry-lineage-enterprise}

Enterprise observability stacks produce high-volume span data that is structurally insufficient for
process mining. The gap is not a tooling deficiency: pm4py, \texttt{wasm4pm}, and related
conformance engines are capable of replaying process models against event logs with measurable
fitness and precision. The gap is a schema deficiency. Standard OTel spans carry no case
identifier that maps to a process instance, no Petri net token state, no lifecycle phase
constrained to a lawful transition set, and no witness seal that can be independently verified.

EXP-001 addresses the enterprise process gap at its root: the feedstock layer. By defining
\texttt{process.pi.instance\_id} as a required attribute (the case identifier in OCEL terms),
the registry ensures that every admissible span can be grouped into a process instance. By requiring
\texttt{process.pi.activity.name} to carry a Petri net transition label rather than an
application-layer string, the registry ensures that the span's activity maps to a declared process
model node. Without these constraints, process mining against enterprise OTel data produces
spaghetti models with low fitness and unmeasurable precision --- the diagnostic signature of
a feedstock schema gap.

\subsection{Relationship to the Chatman Equation}
\label{sec:otel-weaver-exp-001-custom-pi-registry-lineage-chatman-equation}

The Chatman Equation \(A = \mu(O^{*})\) (AUTHOR THESIS) states that architecture is a function
of the manufacturing process applied to the open ontology surface. In EXP-001's context,
\(\mu\) is the Weaver manufacturing pipeline: it takes a YAML registry as input, resolves it
against the OTel standard, and validates that the defined attributes satisfy the schema contract.
\(O^{*}\) is the process intelligence ontology surface --- the set of process-evidence concepts
(instance identity, activity naming, lifecycle state, witness sealing, token state) that must
appear in the feedstock schema. The architecture \(A\) is the \texttt{process.pi.activity} span
group: eight named attributes organized to carry the process-evidence ontology into the telemetry
layer.

The registry is therefore not a configuration file. It is the manufactured form of the process
intelligence ontology, expressed in a standard-conforming schema language that an external
toolchain can validate without consulting the author's intent. The Chatman Equation predicts that
the quality of the architecture is bounded by the quality of the manufacturing process applied to
the ontology; EXP-001 tests whether the Weaver toolchain constitutes a sufficiently rigorous
\(\mu\) for the process intelligence case.

\subsection{Relationship to Receipts and Replay}
\label{sec:otel-weaver-exp-001-custom-pi-registry-lineage-receipts}

The receipt chain is the mechanism by which process claims become auditable. In the dissertation's
architecture, a receipt is a cryptographic artifact that seals a manufacturing event: it records
what was manufactured, from what inputs, under what pipeline version, at what time. EXP-001's
receipt chain is declared in \texttt{otel-weaver-source.manifest.yaml}, which specifies
\texttt{receipt.format: cryptographic-chain}, \texttt{algorithm: blake3}, and
\texttt{store\_location: receipts/otel-weaver-source-receipt.json}.

The relationship to replay is direct: a span that carries \texttt{process.pi.witness.hash} and
\texttt{process.pi.token.state\_before}/\texttt{state\_after} is a span that can be independently
verified and then projected into an OCEL event log for pm4py replay. The schema attributes are
not decorative; they are the fields that the \texttt{OtelSpanToOcelEventProjection} zero-sized
witness type (defined in EXP-003's type system) uses to perform the OTel-to-OCEL translation
without information loss exceeding the \texttt{LossPolicy} threshold. The receipt seals the
manufacturing act; the schema attributes enable the replay act; the two together constitute the
evidence chain that the dissertation's conformance architecture requires.
